% cau-hoi.tex: gói 2 bộ Supply Chain (scm). Module 3 String & Date (9 câu), Module 4 Aggregation (9 câu), case c3, c4.
% SQL nằm trong solution.sql, nối với câu bằng mã câu. Biên dịch: XeLaTeX. Hướng dẫn: HUONG_DAN.md.
\documentclass[11pt,a4paper]{article}
\usepackage{aio-sql-cauhoi}
% Vá XeLaTeX (cong-cu/va-xelatex.tex), không sửa .sty.
\setmonofont{NotoSansMono}[Extension=.ttf, UprightFont=*-Regular, BoldFont=*-Bold, Scale=0.86]
\renewcommand{\erdkhoa}[1]{\node[font=\small, color=cqMuted, anchor=north west, align=left, text width=14cm] at ([yshift=-1mm]current bounding box.south west) {\textbf{Khoá:} #1};}

% ================================================================ KHAI BÁO DÙNG CHUNG
% ---------------------------------------------------------------- bộ dữ liệu
\khaibaodulieu{scm}{Supply Chain · chuỗi bán lẻ và giao hàng chặng cuối}
  {Chuỗi bán lẻ 54 cửa hàng, 3 trung tâm phân phối, 66 nhà cung cấp: lượng bán theo ngày (tháng 7/2017), tồn kho theo kỳ (2017-01 tới 2017-08), đơn mua (2013--2017), đơn giao chặng cuối (2015--2017).}

% ---------------------------------------------------------------- mảng nghiệp vụ
\khaibaodomain{sales}{Bán hàng}{Doanh thu, lượng bán, cửa hàng, ngành hàng.}
\khaibaodomain{marketing}{Marketing}{Khuyến mãi, ngày cao điểm, voucher.}
\khaibaodomain{customer}{Khách hàng}{Lượt giao dịch tại cửa hàng, đơn online theo khách và điểm giao.}
\khaibaodomain{operations}{Vận hành}{Kho, tồn kho, giao hàng, nhà cung cấp, nhân sự kho.}
\khaibaodomain{finance}{Tài chính}{Giá vốn, biên lợi nhuận gộp, vốn tồn kho.}
\khaibaodomain{after_sales}{Hậu mãi}{Hàng lỗi, lô trượt, huỷ đơn.}

% ---------------------------------------------------------------- khái niệm
% \khaibaokhainiem{ma}{Tên}{Định nghĩa}{Công thức}{Trong bộ dữ liệu}
\khaibaokhainiem{ky-va-thang-du}{Kỳ và tháng đủ}
  {Kỳ là đơn vị thời gian của số liệu tồn kho, mỗi kỳ là một tháng. Tháng đủ là tháng có dữ liệu tới ngày cuối tháng.}
  {}
  {\code{inventory_snapshots.ym} là chuỗi dạng \code{'2017-07'}, không phải kiểu ngày. Kỳ 2017-08 chỉ có dữ liệu nửa tháng, nên không dùng khi so các tháng với nhau.}
\khaibaokhainiem{luong-ban-va-nhu-cau}{Lượng bán trong kỳ và nhu cầu}
  {Nhu cầu là lượng khách muốn mua; lượng bán là lượng đã ghi nhận bán ra. Khi hết hàng, lượng bán nhỏ hơn nhu cầu.}
  {}
  {\code{inventory_snapshots.sold} là lượng bán trong kỳ, \code{received} là lượng nhập trong kỳ, \code{thieu_hut} là phần bán vượt lượng hàng sẵn có của tồn kho mô phỏng. Doanh thu, giá vốn tính trên \code{sold}, không trừ \code{thieu_hut}.}
\khaibaokhainiem{von-ton-kho}{Vốn tồn kho}
  {Vốn tồn kho là giá trị hàng còn trong kho cuối kỳ, tính theo giá vốn.}
  {tồn cuối kỳ × giá vốn đơn vị}
  {có sẵn ở \code{inventory_snapshots.inventory_value}, bằng \code{on_hand_end} nhân \code{item_costs.unit_cost}.}
\khaibaokhainiem{bien-ln-gop}{Biên lợi nhuận gộp}
  {Biên lợi nhuận gộp là phần lợi nhuận gộp trên mỗi đồng doanh thu. Lợi nhuận gộp bằng doanh thu trừ giá vốn hàng bán.}
  {(doanh thu - giá vốn hàng bán) / doanh thu}
  {giá cố định ở \code{item_costs} (\code{unit_price}, \code{unit_cost}); lượng bán trong kỳ là \code{inventory_snapshots.sold}. Biên của một nhóm tính lại từ tổng doanh thu và tổng giá vốn, không lấy trung bình biên của từng mã.}
\khaibaokhainiem{lead-time-ncc}{Lead time của nhà cung cấp}
  {Lead time là khoảng thời gian từ khi phát đơn mua tới khi hàng về kho.}
  {ngày nhận hàng - ngày đặt}
  {\code{purchase_orders}: \code{received_date} trừ \code{po_date} (hai cột kiểu ngày); ngày hẹn giao ở \code{promised_date}. Đơn chưa về có \code{received_date} trống.}
\khaibaokhainiem{otd-ncc}{Tỷ lệ giao đúng hạn của nhà cung cấp (OTD)}
  {Tỷ lệ giao đúng hạn là phần đơn mua về không muộn hơn ngày hẹn trong số đơn đã về.}
  {số đơn có received\_date ≤ promised\_date / số đơn đã về}
  {\code{purchase_orders}; đơn chưa về có \code{received_date} trống và \code{inspection_result = 'Pending'}, không vào mẫu số.}
\khaibaokhainiem{pheu-don-mua}{Phễu đơn mua}
  {Phễu đơn mua theo dõi một tập đơn mua qua các bước nối tiếp: đặt hàng, về đúng hạn, về đủ hàng, đạt kiểm định. Mỗi bước chỉ giữ những đơn đã qua mọi bước trước nó.}
  {tỷ lệ qua một bước = số đơn qua bước đó / số đơn qua bước liền trước}
  {\code{purchase_orders}: về đúng hạn khi \code{received_date} không muộn hơn \code{promised_date}; về đủ hàng khi \code{qty_received} bằng \code{qty_ordered}; đạt kiểm định khi \code{inspection_result = 'Pass'}. DC của đơn mua là DC mà nhà cung cấp giao về (\code{suppliers.warehouse_id}).}
\khaibaokhainiem{lo-dat-lo-truot}{Lô đạt, lô trượt}
  {Kiểm định khi nhận hàng xếp mỗi đơn mua đã về thành lô đạt hoặc lô trượt theo ngưỡng hàng lỗi cho phép (AQL).}
  {tỷ lệ lô trượt = số lô trượt / số lô đã về}
  {\code{purchase_orders.inspection_result}: \code{'Pass'} (lô đạt), \code{'Fail'} (lô trượt), \code{'Pending'} (đơn chưa về).}
\khaibaokhainiem{otif-chuyen-giao}{Giao đúng hạn của chuyến giao}
  {Chuyến giao đúng hạn là chuyến từ DC tới cửa hàng có số ngày thực tế không vượt số ngày cam kết. OTIF đòi hỏi thêm giao đủ số lượng.}
  {tỷ lệ giao đúng hạn = số chuyến đúng hạn / tổng số chuyến}
  {\code{shipments}: \code{on_time} bằng 1 khi \code{actual_days} không vượt \code{promised_days}; \code{week} là chuỗi \code{'YYYY-MM-DD/YYYY-MM-DD'} (thứ Hai tới Chủ nhật của tuần giao).}
\khaibaokhainiem{giao-dung-han-don}{Tỷ lệ giao đúng hạn của đơn chặng cuối}
  {Tỷ lệ giao đúng hạn là phần đơn được giao không muộn hơn mức cam kết.}
  {số đơn không huỷ có on\_time = 1 / số đơn không huỷ}
  {\code{orders.on_time} bằng 1 khi \code{lead_time_total_h} không vượt \code{promised_h}. Đơn huỷ (\code{is_cancelled} = 1) vẫn có giá trị ở các cột thời gian, nên phải loại trước khi tính.}
\khaibaokhainiem{lead-time-don}{Lead time của đơn và bốn thành phần}
  {Lead time của đơn là thời gian từ lúc tạo đơn tới lúc giao thành công. Bốn thành phần là t1 (tạo đơn tới lấy hàng), t2 (lấy hàng tới xuất kho), middle mile (t3a di chuyển cộng t3b phân loại) và t4 (kho tới giao thành công).}
  {lead time = delivered\_at - created\_at (giờ)}
  {\code{orders.lead_time_total_h} (giờ, có phần lẻ); \code{created_at}, \code{delivered_at} là mốc thời gian tới phần triệu giây; năm cột chặng \code{t1_tao_don_pick_h}, \code{t2_pick_xuat_h}, \code{t3a_di_chuyen_h}, \code{t3b_sorting_h}, \code{t4_kho_giao_h}. Kho xuất ở \code{origin_hub_id}, kho đích ở \code{dest_hub_id}, cùng nối \code{sorting_hubs}.}
\khaibaokhainiem{gmv}{GMV (tổng giá trị đơn hoàn tất)}
  {GMV là tổng giá trị các đơn hoàn tất, không tính đơn huỷ.}
  {tổng order\_value\_usd của đơn không huỷ}
  {\code{orders.order_value_usd} (USD); đơn huỷ có \code{is_cancelled} = 1.}
\khaibaokhainiem{ty-le-huy}{Tỷ lệ huỷ}
  {Tỷ lệ huỷ là phần đơn bị huỷ trong tổng số đơn được tạo.}
  {số đơn huỷ / tổng số đơn tạo}
  {\code{orders.is_cancelled} bằng 1 là đơn huỷ; \code{order_date} là ngày tạo đơn, \code{created_at} là thời điểm tạo đơn.}
\khaibaokhainiem{gio-chot-don}{Giờ chốt đơn (cut-off time)}
  {Giờ chốt đơn là mốc giờ trong ngày; đơn tạo từ mốc này trở đi được xử lý vào ngày hôm sau.}
  {}
  {giờ tạo đơn ở \code{orders.created_at}, giờ lấy hàng ở \code{picked_at}; giờ chốt đơn hiện nay là 17 giờ 00. Đơn tạo từ 17 giờ trở đi chờ tới sáng hôm sau mới được lấy hàng. Cột \code{t1_tao_don_pick_h} chỉ ghi thời gian xử lý, không gồm phần chờ qua đêm.}
\khaibaokhainiem{ngay-cao-diem}{Ngày cao điểm, ngày nền và ngày thường}
  {Ngày cao điểm là ngày khuyến mãi hoặc ngày lễ trong lịch sự kiện, khi lượng đơn tăng mạnh. Ngày thường là ngày không có trong lịch sự kiện. Ngày nền là ngày nằm ngoài cửa sổ hai ngày trước và sau mỗi ngày cao điểm.}
  {}
  {\code{event_calendar}: mỗi dòng là một ngày sự kiện; \code{tier} là quy mô (\code{'nho'}, \code{'vua'}, \code{'lon'}, \code{'le'}); \code{demand_multiplier} là hệ số nhu cầu.}
\khaibaokhainiem{luong-ban-ngay}{Lượng bán theo ngày}
  {Lượng bán là số đơn vị hàng đã ghi nhận bán ra. Dòng trả hàng được trừ vào tổng.}
  {tổng unit\_sales}
  {bảng \code{sales}: mỗi dòng là một mã hàng tại một cửa hàng trong một ngày; \code{unit_sales} âm là dòng trả hàng. Bộ dữ liệu chỉ có tháng 7/2017; ngành hàng của mã lấy từ \code{items.family}.}
\khaibaokhainiem{khuyen-mai}{Khuyến mãi và mức tăng lượng bán (lift)}
  {Khuyến mãi là giảm giá hoặc trưng bày đặc biệt cho một mã hàng trong một ngày. Mức tăng lượng bán so sánh lượng bán bình quân khi có và khi không có khuyến mãi.}
  {lượng bán bình quân khi có khuyến mãi / lượng bán bình quân khi không khuyến mãi}
  {\code{sales.onpromotion} bằng 1 là dòng có khuyến mãi. Mã được chọn khuyến mãi thường đã bán chạy sẵn, nên chênh lệch thô chỉ cho thấy tương quan.}
\khaibaokhainiem{luot-giao-dich}{Lượt giao dịch}
  {Lượt giao dịch là số lần thanh toán tại quầy của một cửa hàng trong ngày, dùng để đo mức đông khách.}
  {}
  {bảng \code{transactions}: mỗi dòng là một cửa hàng trong một ngày (\code{date}, \code{store_nbr}, \code{transactions}).}
\khaibaokhainiem{voucher}{Voucher (phiếu giảm giá)}
  {Voucher là mã giảm giá áp vào đơn online. Khách có thể tự áp mã hoặc được hệ thống tự động áp mã khi thanh toán.}
  {tỷ lệ đơn dùng voucher = số đơn có voucher / tổng số đơn}
  {\code{orders.voucher_applied} bằng 1 là đơn có voucher; \code{checkout_group} là nhóm khách tự áp mã (\code{'manual'}) hoặc nhóm hệ thống tự động áp mã (\code{'auto_apply'}); đơn của nhóm \code{'auto_apply'} luôn có \code{voucher_applied} = 1.}

\begin{document}
{\LARGE\bfseries Gói 2 · Supply Chain · String \& Date, Aggregation\par}
\medskip
{\color{cqMuted}Bộ dữ liệu \texttt{scm} · Module 3 (9 câu), Module 4 (9 câu), 2 case · chuẩn \texttt{aio-sql-cauhoi} v0.4\par}
\tableofcontents


% ################################################################ MODULE 3 · STRING & DATE

% ================================================================ S1 · write · Cơ bản
\begin{cauhoi}[ma=scm-s1-don-mua-ve-tre-2016, dulieu=scm, domain=operations,
  hinhthuc=write, cap=basic, dang=A1, kynang={DATEDIFF, ORDER BY, LIMIT}, bay={dao-thu-tu-datediff},
  tacgia={Agent gói 2}, trangthai=review, bokiemtra={public, hidden-1}, thutu=co]
\tieude{Đơn mua nào về trễ hạn nhất năm 2016?}
\boicanh[Chị Hằng · Trưởng phòng mua hàng]{Tuần sau chị họp với các nhà cung cấp hay giao trễ. Em lấy giúp chị 20 đơn mua năm 2016 về muộn nhất so với ngày hẹn.}
\begin{yeucau}
  \item Xét các đơn mua có \code{po_date} trong năm 2016 và về sau ngày hẹn giao (\code{received_date} lớn hơn \code{promised_date}).
  \item Tính \code{so_ngay_tre} là số ngày từ \code{promised_date} tới \code{received_date} (số dương).
  \item Trả về \code{po_id}, \code{supplier_id}, \code{promised_date}, \code{received_date}, \code{so_ngay_tre}.
  \item Sắp \code{so_ngay_tre} giảm dần; bằng nhau thì \code{po_id} nhỏ hơn đứng trước. Chỉ lấy 20 dòng đầu.
\end{yeucau}
\begin{erd}
  \erdbang{purchase_orders}{0,0}{po_id:PK, supplier_id:FK, po_date*=2016-03-31, promised_date*=2016-04-07, received_date*=2016-04-12}
\end{erd}
\begin{nghiepvu}{Chọn nhà cung cấp cần nhắc nhở về lịch giao hàng.}
  \quytrinh{Vòng đời đơn mua}{đặt hàng, hàng về, kiểm định}{chưa về}
  \khainiem{lead-time-ncc}
  \luuy{Số ngày trễ tính từ ngày hẹn giao, không tính từ ngày đặt. Đơn về đúng hoặc sớm hơn ngày hẹn không có ngày trễ.}
\end{nghiepvu}
\goiy{\code{DATEDIFF(a, b)} trả về a trừ b. Ngày nào phải đứng trước để số ngày trễ là số dương?}
\ghichuthietke{Cơ bản: 1 bảng, giá trị mẫu cho thấy ba cột ngày đã là kiểu DATE. Bài sai đảo đối số của DATEDIFF ra số âm nên sắp giảm dần lấy nhầm đơn trễ ít nhất (rớt ở public).}
\end{cauhoi}

% ================================================================ S2 · write · Cơ bản
\begin{cauhoi}[ma=scm-s2-ngay-cao-diem-2016, dulieu=scm, domain=marketing,
  hinhthuc=write, cap=basic, dang=A1, kynang={DATE_FORMAT, DAYOFWEEK}, bay={sap-theo-chuoi-ngay},
  tacgia={Agent gói 2}, trangthai=review, bokiemtra={public, hidden-1}, thutu=co]
\tieude{Ngày cao điểm năm 2016 rơi vào thứ mấy?}
\boicanh[Chị Ngọc · Trưởng nhóm marketing]{Chị đang lên lịch truyền thông cho mùa cao điểm năm nay. Em liệt kê giúp chị các ngày cao điểm năm ngoái (2016) và mỗi ngày rơi vào thứ mấy.}
\begin{yeucau}
  \item Xét các ngày cao điểm có \code{date} trong năm 2016.
  \item Trả về \code{event_name}; \code{ngay} là chuỗi dạng \code{'dd/mm/yyyy'} (ví dụ \code{'09/09/2016'}); \code{thu} là thứ trong tuần dạng số, 1 là Chủ nhật, 2 là thứ Hai, …, 7 là thứ Bảy; \code{tier}.
  \item Sắp theo ngày diễn ra, sớm trước.
\end{yeucau}
\begin{erd}
  \erdbang{event_calendar}{0,0}{date*=2016-11-11, event_name*='Sale 11.11', tier*='lon', demand_multiplier}
\end{erd}
\begin{nghiepvu}{Xếp lịch truyền thông trước các ngày cao điểm.}
  \khainiem{ngay-cao-diem}
  \luuy{Cùng một sự kiện, mỗi năm rơi vào một thứ khác nhau trong tuần. Ngày cao điểm trùng thứ Bảy hoặc Chủ nhật có lượng đơn khác ngày cao điểm rơi vào ngày làm việc.}
\end{nghiepvu}
\goiy{Chuỗi \code{'dd/mm/yyyy'} sắp theo chữ cái sẽ không theo thứ tự thời gian. Sắp theo cột nào thì đúng?}
\goiy{\code{DATE_FORMAT(`date`, '\%d/\%m/\%Y')} và \code{DAYOFWEEK(`date`)}.}
\ghichuthietke{Cơ bản: 1 bảng có giá trị mẫu. Hai bài sai rớt ở public: sắp theo chuỗi \code{ngay}; dùng \code{WEEKDAY} (0 là thứ Hai) thay \code{DAYOFWEEK}.}
\end{cauhoi}

% ================================================================ S3 · debug · Trung bình
\begin{cauhoi}[ma=scm-s3-huy-kho-ve-tinh-18-thang, dulieu=scm, domain=after_sales,
  hinhthuc=debug, cap=medium, dang=B1, kynang={DATE_FORMAT, GROUP BY, JOIN}, bay={gom-theo-thang-mat-nam},
  tacgia={Agent gói 2}, trangthai=review, bokiemtra={public, hidden-1}, thutu=co]
\tieude{Báo cáo tỷ lệ huỷ 18 tháng chỉ còn 12 dòng}
\boicanh[Chị Hà · Trưởng vận hành]{Chị cần xu hướng huỷ đơn ở kho vệ tinh từ đầu năm 2016 tới hết tháng 6/2017 để họp quý. Báo cáo em gửi chỉ có 12 tháng của năm 2016.}
\begin{yeucau}
  \item Truy vấn dưới đây tính số đơn, số đơn huỷ và tỷ lệ huỷ theo tháng đặt đơn, cho các đơn có kho đích là kho vệ tinh, đặt từ tháng 1/2016 tới hết tháng 6/2017. Kết quả chỉ có 12 dòng.
  \item Sửa để mỗi tháng trong 18 tháng là một dòng; giữ nguyên cột \code{thang} (chuỗi dạng \code{'2016-01'}), \code{so_don}, \code{so_huy}, \code{ty_le_huy_pct} (phần trăm, làm tròn 2 chữ số).
  \item Sắp \code{thang} tăng dần.
\end{yeucau}
\begin{erd}
  \erdbang{orders}{0,0}{order_id:PK, order_date*, dest_hub_id:FK, dest_point_id:FK, is_cancelled*}
  \erdbang{sorting_hubs}{1,0}{hub_id:PK, role*, city}
  \erdbang{delivery_points}{1,1}{point_id:PK, city}
  \erdnoi[kieu=n-1]{orders.dest_hub_id}{sorting_hubs.hub_id}
  \erdnoi[kieu=n-1]{orders.dest_point_id}{delivery_points.point_id}
\end{erd}
\begin{nghiepvu}{Theo dõi xu hướng huỷ đơn ở kho vệ tinh qua từng tháng để rà lại quy trình xử lý đơn.}
  \khainiem{ty-le-huy}
  \luuy{Tỷ lệ huỷ của một tháng tính trên mọi đơn đặt trong tháng đó, kể cả đơn huỷ.}
\end{nghiepvu}
\goiy{Đếm riêng số đơn tháng 1/2016 rồi so với dòng \code{'2016-01'} của báo cáo. Phần chênh là đơn của tháng nào?}
\goiy{Gom nhóm theo \code{DATE_FORMAT(o.order_date, '\%Y-\%m')}.}
\ghichuthietke{Trung bình (sửa v4, viết lại từ câu Cơ bản \code{scm-s3-lo-truot-thang-3-ra-0}, câu cũ trùng ô với \code{scm-r7}). Lỗi khó thấy: starter gom theo \code{MONTH(order_date)} nên tháng 1/2016 và tháng 1/2017 dồn vào một nhóm; nhãn \code{DATE_FORMAT(MIN(order_date), ...)} lại hiện tháng của năm 2016 nên bảng trông như đúng. 5 trục: 2 bảng dùng (Trung bình); 2 tầng, nối rồi gom (Dễ); tỷ lệ huỷ định nghĩa ở tab Nghiệp vụ (Trung bình); 1 bẫy khó thấy (Trung bình); công thức hai thành phần (Trung bình) → Trung bình 4/5. ERD 3 bảng chỉ ký hiệu quan hệ, \code{delivery_points} là bảng nhiễu. Bài sai: nhãn và nhóm theo \code{'\%Y-\%M'} (tên tháng tiếng Anh, sắp theo chữ); gom theo năm và tháng nhưng ghép \code{CONCAT(YEAR, '-', MONTH)} mất số 0; cả hai rớt ở public.}
\end{cauhoi}

% ================================================================ S4 · write · Khó
\begin{cauhoi}[ma=scm-s4-lui-gio-chot-don, dulieu=scm, domain=operations,
  hinhthuc=write, cap=hard, dang=A2, kynang={HOUR, TIMESTAMPDIFF, CASE WHEN, CTE}, bay={khung-gio-chot, tinh-ca-don-huy, gio-lam-tron, quen-can-tren},
  tacgia={Agent gói 2}, trangthai=review, bokiemtra={public, hidden-1}, thutu=co]
\tieude{Lead time giảm bao nhiêu nếu lùi giờ chốt đơn?}
\boicanh[Anh Dũng · Điều phối vận hành]{Các kho đề xuất lùi giờ chốt đơn thêm một tiếng, sang 18 giờ. Trước khi trình ban giám đốc, anh cần ước tính lead time bình quân của từng kho xuất sẽ còn bao nhiêu.}
\begin{yeucau}
  \item Xét các đơn không huỷ, đặt từ ngày 1/7/2017 tới hết tháng 8/2017.
  \item Giả định khi lùi giờ chốt sang 18 giờ, đơn tạo trong khung 17 giờ (từ 17:00 tới trước 18:00) được lấy hàng ngay trong ngày nên không còn phần chờ qua đêm; lead time của các đơn khác giữ nguyên.
  \item Tính cho từng kho xuất số đơn, số đơn tạo trong khung 17 giờ, lead time bình quân hiện nay và lead time bình quân sau khi lùi giờ chốt, theo tab Nghiệp vụ.
  \item Trả về \code{origin_hub_id}, \code{so_don}, \code{so_don_khung_17h}, \code{lead_time_tb}, \code{lead_time_tb_moi} (hai cột cuối tính bằng giờ, làm tròn 2 chữ số).
  \item Sắp \code{origin_hub_id} tăng dần.
\end{yeucau}
\begin{erd}
  \erdbang{orders}{0,0}{order_id:PK, origin_hub_id:FK, dest_hub_id:FK, dest_point_id:FK}
  \erdbang{sorting_hubs}{1,0}{hub_id:PK}
  \erdbang{warehouse_daily_ops}{2,0}{hub_id:FK, date}
  \erdbang{staffing_plan}{2,1}{hub_id:FK, date}
  \erdbang{delivery_points}{0,1}{point_id:PK}
  \erdkhoa{Mọi cột cần dùng nằm trong \code{orders}; kho xuất là \code{origin_hub_id}. \code{sorting_hubs}, \code{warehouse_daily_ops}, \code{staffing_plan}, \code{delivery_points} không cần.}
\end{erd}
\begin{nghiepvu}{Quyết định lùi giờ chốt đơn hay giữ giờ chốt và tăng ca lấy hàng buổi sáng.}
  \quytrinh{Hành trình đơn chặng cuối}{tạo đơn, lấy hàng, xuất kho, phân loại, giao}{huỷ}
  \khainiem{gio-chot-don}
  \khainiem{lead-time-don}
  \luuy{Lead time bình quân của một kho tính trên mọi đơn không huỷ của kho đó, kể cả đơn tạo ngoài khung 17 giờ. Đơn huỷ không được giao nên không có lead time.}
\end{nghiepvu}
\goiy{Phần chờ qua đêm của một đơn là khoảng từ \code{created_at} tới \code{picked_at} trừ đi \code{t1_tao_don_pick_h}. Tính khoảng này theo giây rồi đổi ra giờ để không bị làm tròn.}
\goiy{Lead time mới của kho bằng \code{(SUM(lead_time_total_h) - SUM(gio_cho_khung_17h)) / COUNT(*)}.}
\ghichuthietke{Khó (sửa v4, viết lại từ câu Cơ bản \code{scm-s4-don-sau-gio-chot-thang-7}). Giờ chốt 17 giờ theo sách dự án (phần kiểm tra năm cột chặng): mọi đơn tạo từ 17 giờ trở đi chờ qua đêm, \code{t1} không ghi phần chờ; khái niệm \code{gio-chot-don} sửa từ 18 giờ thành 17 giờ cho khớp sách và dữ liệu. 5 trục: 1 bảng dùng (chưa đạt, như mẫu o4); 3 tầng, giờ chờ từng đơn rồi gom theo kho rồi tính lead time mới (Trung bình); đề nêu giả định kinh doanh, học viên tự suy ra giờ chờ từ khái niệm (Khó); 4 bẫy, 1 chỉ lộ ở hidden-1 (Khó); chỉ số nhiều điều kiện: đơn không huỷ, khung giờ, giờ chờ trừ thời gian xử lý, chia trên mọi đơn (Khó) → Khó 3/5. ERD 5 bảng chỉ khoá, 4 bảng nhiễu (lời giải dùng 1 bảng), không vẽ đường nối. Bài sai: tính cả đơn huỷ; coi mọi đơn từ 17 giờ trở đi đều hết chờ; đổi giờ chờ bằng \code{TIMESTAMPDIFF(HOUR, ...)} (làm tròn xuống giờ chẵn); chia phần giờ tiết kiệm cho số đơn khung 17 giờ; bốn bài rớt ở public. Bài sai thiếu cận trên rớt ở hidden-1 (đơn \code{ORD2_01033328} ngày 2017-09-01).}
\end{cauhoi}

% ================================================================ S5 · write · Dễ
\begin{cauhoi}[ma=scm-s5-ban-theo-thu-cua-hang-44, dulieu=scm, domain=sales,
  hinhthuc=write, cap=easy, dang=B1, kynang={DAYOFWEEK, GROUP BY, JOIN}, bay={quen-can-tren},
  tacgia={Agent gói 2}, trangthai=review, bokiemtra={public, hidden-1}, thutu=co]
\tieude{Cửa hàng 44 bán BEVERAGES nhiều nhất vào thứ nào?}
\boicanh[Chị Thảo · Quản lý cửa hàng]{Chị muốn xếp lịch bổ sung đồ uống lên kệ theo nhịp bán trong tuần. Em xem giúp chị tháng 7/2017 cửa hàng mình (cửa hàng 44) bán ngành BEVERAGES nhiều nhất vào thứ nào.}
\begin{yeucau}
  \item Xét các dòng bán của cửa hàng 44, ngành hàng \code{'BEVERAGES'}, từ ngày 1/7/2017 tới hết ngày 31/7/2017.
  \item Tính tổng lượng bán của từng thứ trong tuần cho cả tháng (không chia cho số ngày), theo tab Nghiệp vụ.
  \item Trả về \code{thu} (1 là Chủ nhật, …, 7 là thứ Bảy), \code{luong_ban} (đơn vị, làm tròn 0 chữ số).
  \item Sắp \code{thu} tăng dần.
\end{yeucau}
\begin{erd}
  \erdbang{sales}{0,0}{date*, store_nbr*, item_nbr:FK, unit_sales*, onpromotion}
  \erdbang{items}{1,0}{item_nbr:PK, family*, class, perishable}
  \erdnoi[kieu=n-1, tren={\code{sales.item_nbr = items.item_nbr}}, duoi={nhiều dòng bán, một mã hàng}]{sales.item_nbr}{items.item_nbr}
\end{erd}
\begin{nghiepvu}{Xếp lịch bổ sung hàng lên kệ theo nhịp bán trong tuần.}
  \khainiem{luong-ban-ngay}
  \luuy{Tháng 7/2017 bắt đầu vào thứ Bảy, nên thứ Bảy, Chủ nhật và thứ Hai có 5 ngày, các thứ khác có 4 ngày. Khi so các thứ bằng tổng của cả tháng, cần nhớ ba thứ này có thêm một ngày.}
\end{nghiepvu}
\goiy{Ngành hàng nằm ở \code{items.family}; nối \code{sales} với \code{items} theo \code{item_nbr}.}
\goiy{Gom theo \code{DAYOFWEEK(`date`)}; lọc ngày có cả cận dưới và cận trên.}
\ghichuthietke{Dễ: 2 bảng, điều kiện nối ghi trên đường nối. Bài sai thiếu cận trên rớt ở hidden-1 (thêm dòng BEVERAGES của cửa hàng 44 ngày 2017-08-01, thứ Ba); bài sai dùng \code{WEEKDAY} rớt ở public. Lọc một cửa hàng để truy vấn dùng chỉ mục \code{idx_sales_store_item}.}
\end{cauhoi}

% ================================================================ S6 · write · Dễ
\begin{cauhoi}[ma=scm-s6-chuyen-giao-tre-quy-2, dulieu=scm, domain=operations,
  hinhthuc=write, cap=easy, dang=A1, kynang={SUBSTRING_INDEX, JOIN}, bay={lay-sai-phan-chuoi},
  tacgia={Agent gói 2}, trangthai=review, bokiemtra={public, hidden-1}, thutu=co]
\tieude{Chuyến giao nào trễ hạn trong quý 2/2017?}
\boicanh[Anh Dũng · Điều phối vận hành]{Anh cần danh sách các chuyến giao trễ từ trung tâm phân phối (DC) tới cửa hàng trong quý 2/2017 để làm việc với đội xe. Em lấy giúp anh kèm thành phố của cửa hàng.}
\begin{yeucau}
  \item Xét các chuyến giao có \code{on_time} bằng 0 và ngày đầu tuần giao từ \code{'2017-04-01'} tới hết \code{'2017-06-30'}.
  \item Ngày đầu tuần giao \code{tuan_bat_dau} là phần trước dấu \code{/} của cột \code{week}, giữ dạng chuỗi \code{'YYYY-MM-DD'}.
  \item Trả về \code{shipment_id}, \code{tuan_bat_dau}, \code{store_nbr}, \code{city} (thành phố của cửa hàng), \code{so_ngay_tre} (\code{actual_days} trừ \code{promised_days}, làm tròn 2 chữ số).
  \item Sắp \code{tuan_bat_dau} tăng dần, rồi \code{shipment_id} tăng dần.
\end{yeucau}
\begin{erd}
  \erdbang{shipments}{0,0}{shipment_id:PK, store_nbr:FK, week*, promised_days*, actual_days*, on_time*}
  \erdbang{stores}{1,0}{store_nbr:PK, city*, state, type}
  \erdnoi[kieu=n-1, tren={\code{shipments.store_nbr = stores.store_nbr}}, duoi={nhiều chuyến, một cửa hàng}]{shipments.store_nbr}{stores.store_nbr}
\end{erd}
\begin{nghiepvu}{Chọn tuyến và đội xe cần nhắc nhở về thời gian giao.}
  \khainiem{otif-chuyen-giao}
  \luuy{Một tuần giao có thể bắt đầu ở quý trước và kết thúc ở quý sau. Tuần giao được tính vào quý của ngày đầu tuần.}
\end{nghiepvu}
\goiy{Cột \code{week} có dạng \code{'2017-04-03/2017-04-09'}. Phần nào là ngày đầu tuần?}
\goiy{\code{SUBSTRING_INDEX(week, '/', 1)} lấy phần trước dấu \code{/} đầu tiên.}
\ghichuthietke{Dễ: 2 bảng có điều kiện nối. Bài sai lấy phần sau dấu \code{/} (ngày cuối tuần) lệch một tuần ở hai đầu quý, rớt ở public.}
\end{cauhoi}

% ================================================================ S7 · debug · Dễ
\begin{cauhoi}[ma=scm-s7-gia-tri-nhan-sai-thu-tu, dulieu=scm, domain=finance,
  hinhthuc=debug, cap=easy, dang=B1, kynang={DATE_FORMAT, GROUP BY}, bay={sap-ten-thang},
  tacgia={Agent gói 2}, trangthai=review, bokiemtra={public, hidden-1}, thutu=co]
\tieude{Báo cáo giá trị hàng nhận xếp sai thứ tự tháng}
\boicanh[Anh Minh · Kế toán quản trị]{Báo cáo giá trị hàng nhận năm 2016 của em đưa tháng 10, 11, 12 lên ngay sau tháng 1. Em sửa giúp anh để các tháng xếp đúng thứ tự từ tháng 1 tới tháng 12.}
\begin{yeucau}
  \item Truy vấn dưới đây tính giá trị hàng nhận theo giá vốn cho từng tháng nhận hàng năm 2016, nhưng các tháng xếp sai thứ tự.
  \item Sửa để \code{thang} có dạng \code{'2016-01'} và các dòng xếp đúng thứ tự thời gian.
  \item Giữ nguyên cột \code{thang}, \code{gia_tri_nhan} (làm tròn 0 chữ số).
\end{yeucau}
\begin{erd}
  \erdbang{purchase_orders}{0,0}{po_id:PK, item_nbr:FK, received_date*, qty_received*}
  \erdbang{item_costs}{1,0}{item_nbr:PK, family, unit_price, unit_cost*}
  \erdnoi[kieu=n-1, tren={\code{purchase_orders.item_nbr = item_costs.item_nbr}}]{purchase_orders.item_nbr}{item_costs.item_nbr}
\end{erd}
\begin{nghiepvu}{Theo dõi giá trị hàng nhập kho theo tháng để lập kế hoạch dòng tiền mua hàng.}
  \luuy{Giá trị hàng nhận tính theo giá vốn đơn vị cố định ở \code{item_costs}, nhân số lượng thực nhận, không nhân số lượng đặt.}
\end{nghiepvu}
\goiy{Chuỗi tháng sắp theo chữ cái chỉ đúng thứ tự thời gian khi tháng là hai chữ số.}
\goiy{\code{DATE_FORMAT(received_date, '\%Y-\%m')} cho tháng luôn có hai chữ số.}
\ghichuthietke{Dễ: 2 bảng có điều kiện nối. Starter ghép \code{YEAR} và \code{MONTH} bằng \code{CONCAT} nên tháng không có số 0 đứng đầu, \code{'2016-10'} đứng trước \code{'2016-2'} (khác lỗi \code{\%M} của s3). Bài sai giữ \code{CONCAT} nhưng sắp theo ngày nhận sớm nhất: đúng thứ tự mà \code{thang} sai dạng, rớt ở public.}
\end{cauhoi}

% ================================================================ S8 · write · Trung bình
\begin{cauhoi}[ma=scm-s8-von-ton-cuoi-ky-2017, dulieu=scm, domain=finance,
  hinhthuc=write, cap=medium, dang=B1, kynang={CONCAT, STR_TO_DATE, LAST_DAY, GROUP BY}, bay={thang-chua-du, quen-can-tren},
  tacgia={Agent gói 2}, trangthai=review, bokiemtra={public, hidden-1}, thutu=co]
\tieude{Cửa hàng 1 tồn bao nhiêu vốn cuối mỗi kỳ năm 2017?}
\boicanh[Chị Lan · Giám đốc tài chính]{Ngân hàng đang xét khoản vay cải tạo cửa hàng 1. Chị cần số vốn tồn kho cuối từng tháng năm nay của cửa hàng này, tháng nào cũng ghi rõ ngày chốt sổ.}
\begin{yeucau}
  \item Xét các kỳ từ 2017-01 trở đi, chỉ giữ tháng đủ theo tab Nghiệp vụ.
  \item Tính vốn tồn kho cuối kỳ của cửa hàng 1 cho từng kỳ, bằng tổng cột \code{inventory_value} có sẵn (không tính lại từ \code{on_hand_end} và giá vốn).
  \item Trả về \code{ym}; \code{ngay_chot} là ngày cuối tháng của kỳ, kiểu ngày (DATE); \code{von_ton} (làm tròn 0 chữ số).
  \item Sắp \code{ym} tăng dần.
\end{yeucau}
\begin{erd}
  \erdbang{inventory_snapshots}{0,0}{store_nbr, item_nbr:FK, ym*, on_hand_end, inventory_value*}
  \erdbang{item_costs}{1,0}{item_nbr:PK, unit_cost}
  \erdbang{items}{2,0}{item_nbr:PK, family}
  \erdnoi[kieu=n-1]{inventory_snapshots.item_nbr}{item_costs.item_nbr}
  \erdnoi[kieu=1-1]{item_costs.item_nbr}{items.item_nbr}
\end{erd}
\begin{nghiepvu}{Lập báo cáo vốn tồn kho theo tháng cho ngân hàng và ban giám đốc.}
  \khainiem{ky-va-thang-du}
  \khainiem{von-ton-kho}
  \luuy{Kỳ là chuỗi \code{'YYYY-MM'}, không có ngày. Ngày chốt sổ phải dựng từ chuỗi kỳ rồi lấy ngày cuối tháng.}
\end{nghiepvu}
\goiy{Kỳ nào chưa đủ tháng? Điều kiện lọc kỳ cần chặn cả phía trên, không chỉ bỏ riêng một kỳ.}
\goiy{\code{CONCAT(ym, '-01')} cho ngày đầu tháng; \code{STR_TO_DATE} đổi chuỗi thành ngày; \code{LAST_DAY} cho ngày cuối tháng.}
\ghichuthietke{Trung bình: 3 bảng, chỉ ký hiệu quan hệ; vốn tồn đã có sẵn ở \code{inventory_value} nên không cần nối. Bài sai \code{ym <> '2017-08'} rớt ở hidden-1 (có kỳ 2017-09 ở cửa hàng 1); bài sai lấy cả 2017-08 rớt ở public. Sửa V4 (11/11): thu phạm vi từ toàn chuỗi về cửa hàng 1, vì gom cả bảng theo \code{ym} làm MySQL quét theo chỉ mục kỳ (5,5--15 giây); lọc một cửa hàng dùng khoá chính, dưới 0,5 giây. Chọn cửa hàng 1 vì dòng kỳ 2017-09 của hidden-1 nằm ở cửa hàng này, giữ được bẫy chỉ bỏ riêng kỳ 2017-08.}
\end{cauhoi}

% ================================================================ S9 · debug · Khó
\begin{cauhoi}[ma=scm-s9-muc-tang-ngay-cao-diem-2016, dulieu=scm, domain=marketing,
  hinhthuc=debug, cap=hard, dang=C3, kynang={DATE_SUB, DATE_ADD, NOT EXISTS, CTE}, bay={cong-tru-so-vao-ngay, bo-su-kien-ngoai-nam},
  tacgia={Agent gói 2}, trangthai=review, bokiemtra={public, hidden-1}, thutu=co]
\tieude{Mức tăng đơn ngày cao điểm 2016 bị tính thấp}
\boicanh[Chị Ngọc · Trưởng nhóm marketing]{Chị dùng báo cáo này để đặt hệ số nhu cầu cho mùa cao điểm năm nay. Dò tay vài ngày thì chị thấy có ngày sát sự kiện vẫn bị xếp vào ngày nền.}
\begin{yeucau}
  \item Xét các ngày trong năm 2016. Số đơn của một ngày là mọi đơn đặt trong ngày đó, kể cả đơn huỷ.
  \item Mức tăng thực tế của một quy mô sự kiện là số đơn bình quân của các ngày cao điểm thuộc quy mô đó so với số đơn bình quân của ngày nền, đổi ra phần trăm; mức tăng kế hoạch đổi từ hệ số nhu cầu, theo tab Nghiệp vụ.
  \item Truy vấn dưới đây có hai lỗi chồng nhau. Sửa cả hai; giữ nguyên các cột \code{tier}, \code{so_ngay}, \code{muc_tang_ke_hoach_pct}, \code{muc_tang_thuc_te_pct} (phần trăm, làm tròn 2 chữ số) và thứ tự sắp.
\end{yeucau}
\begin{erd}
  \erdbang{orders}{0,0}{order_id:PK, origin_hub_id:FK, dest_hub_id:FK}
  \erdbang{event_calendar}{1,0}{date:PK}
  \erdbang{holidays}{2,0}{date, locale_name}
  \erdbang{sorting_hubs}{0,1}{hub_id:PK}
  \erdbang{warehouse_daily_ops}{1,1}{hub_id:FK, date}
  \erdkhoa{\code{event_calendar} không có khoá chung với \code{orders}; hai bảng gặp nhau qua ngày, \code{orders.order_date} với \code{event_calendar.date}. \code{holidays}, \code{sorting_hubs}, \code{warehouse_daily_ops} không cần.}
\end{erd}
\begin{nghiepvu}{Đặt hệ số nhu cầu cho từng quy mô sự kiện khi lập kế hoạch nhân sự và tồn kho mùa cao điểm.}
  \khainiem{ngay-cao-diem}
  \luuy{Hệ số nhu cầu là số lần số đơn so với ngày nền; mức tăng là phần vượt ngày nền, tính theo phần trăm. Cửa sổ hai ngày tính theo ngày lịch, kể cả khi rơi sang tháng khác hoặc năm khác.}
\end{nghiepvu}
\goiy{Liệt kê các ngày truy vấn đang giữ làm ngày nền rồi soát những ngày cuối tháng 1 và cuối tháng 12/2016. Ngày nào nằm trong hai ngày quanh một ngày cao điểm?}
\goiy{\code{n.ngay BETWEEN DATE_SUB(e.`date`, INTERVAL 2 DAY) AND DATE_ADD(e.`date`, INTERVAL 2 DAY)}, xét mọi ngày cao điểm trong lịch.}
\ghichuthietke{Khó (sửa v4, viết lại từ câu Dễ \code{scm-s9-huy-cuoi-tuan-kho-ve-tinh}). Hai lỗi chồng nhau ở tập ngày nền: (1) \code{e.`date` - 2}, \code{e.`date` + 2} là phép cộng trừ số, ngày 2016-02-02 trừ 2 thành số 20160200 nên ngày 31/1/2016 lọt ra ngoài cửa sổ; (2) chỉ xét sự kiện trong năm 2016 nên bỏ sót cửa sổ của ngày 1/1/2017 (ngày 30 và 31/12/2016). Ba ngày đông đơn bị tính là ngày nền nên mức tăng thực tế thấp đi; sửa một lỗi chưa đủ vì ngày 30, 31/12 vẫn sai. 5 trục: 2 bảng dùng (chưa đạt); 4 tầng, đơn theo ngày, ngày nền, gom theo quy mô, so với nền (Khó); định nghĩa ngày nền ở tab Nghiệp vụ, học viên tự dò ngày sai (Khó); 2 bẫy khó thấy, chưa có bẫy chỉ lộ ở hidden-1 (hidden-1 thêm ngày cao điểm 2016-09-02 nên kết quả khác public); chỉ số nhiều điều kiện (Khó) → Khó 3/5. ERD 5 bảng chỉ khoá, 3 bảng nhiễu, không vẽ đường nối. Bài sai: chỉ sửa phép cộng trừ ngày; chỉ bỏ giới hạn năm của sự kiện; coi mọi ngày không có sự kiện là ngày nền; cả ba rớt ở public.}
\end{cauhoi}


% ################################################################ MODULE 4 · AGGREGATION

% ================================================================ A1 · write · Cơ bản
\begin{cauhoi}[ma=scm-a1-chuyen-giao-theo-dc-2016, dulieu=scm, domain=operations,
  hinhthuc=write, cap=basic, dang=A3, kynang={COUNT, SUM, GROUP BY},
  tacgia={Agent gói 2}, trangthai=review, bokiemtra={public, hidden-1}, thutu=co]
\tieude{Mỗi DC chạy bao nhiêu chuyến giao năm 2016?}
\boicanh[Anh Dũng · Điều phối vận hành]{Cuối năm anh tổng kết khối lượng việc của ba trung tâm phân phối (DC). Em đếm giúp anh năm 2016 mỗi DC chạy bao nhiêu chuyến tới cửa hàng, trong đó bao nhiêu chuyến đúng hạn.}
\begin{yeucau}
  \item Xét các chuyến giao có \code{week} bắt đầu bằng \code{'2016'} (tuần giao bắt đầu trong năm 2016).
  \item Đếm cho từng DC (\code{warehouse_id}) số chuyến và số chuyến đúng hạn (\code{on_time} bằng 1).
  \item Trả về \code{warehouse_id}, \code{so_chuyen}, \code{so_chuyen_dung_han}.
  \item Sắp \code{warehouse_id} tăng dần.
\end{yeucau}
\begin{erd}
  \erdbang{shipments}{0,0}{shipment_id:PK, warehouse_id*='DC_HCM', store_nbr:FK, week*='2016-03-14/2016-03-20', on_time*=1}
\end{erd}
\begin{nghiepvu}{Phân bổ đội xe và nhân sự cho ba DC theo khối lượng chuyến giao.}
  \khainiem{otif-chuyen-giao}
  \luuy{Mỗi dòng của \code{shipments} là một chuyến giao từ một DC tới một cửa hàng trong một tuần.}
\end{nghiepvu}
\goiy{Gom nhóm theo \code{warehouse_id}; \code{on_time} chỉ nhận giá trị 0 hoặc 1.}
\goiy{Lọc \code{week LIKE '2016\%'}.}
\ghichuthietke{Cơ bản (sửa v4, hạ từ câu Dễ \code{scm-a1-dung-han-theo-dc-2016}: bỏ nối \code{warehouses}, bỏ tỷ lệ, ghi rõ giá trị lọc). 5 trục đều ở mức Cơ bản: 1 bảng; 1 tầng, lọc rồi gom theo một cột; đề ghi rõ cột và giá trị lọc, ERD có giá trị mẫu; 0 bẫy; không công thức. Bài sai \code{week LIKE '\%2016\%'} lấy thêm tuần bắt đầu ngày 28/12/2015; bài sai \code{COUNT(on_time)} đếm cả chuyến trễ; cả hai rớt ở public. Hidden-1 bỏ một chuyến và sửa cờ đúng hạn của một phần mười số chuyến nên kết quả khác public.}
\end{cauhoi}

% ================================================================ A2 · debug · Khó
\begin{cauhoi}[ma=scm-a2-dai-chi-so-kho-dich-2017, dulieu=scm, domain=operations,
  hinhthuc=debug, cap=hard, dang=A2, kynang={AVG(CASE WHEN), SUM, GROUP BY}, bay={hai-tap-don, else-0-trong-avg, quen-can-tren},
  tacgia={Agent gói 2}, trangthai=review, bokiemtra={public, hidden-1}, thutu=co]
\tieude{Bảng chỉ số kho đích lệch số đội giao hàng báo}
\boicanh[Chị Hà · Trưởng vận hành]{Bảng chỉ số tháng 7 và tháng 8/2017 theo kho đích ghi lead time cao hơn, tỷ lệ đúng hạn thấp hơn số đội giao hàng báo lên. Chị sắp gửi bản này cho ban giám đốc, em rà lại giúp chị.}
\begin{yeucau}
  \item Bảng chỉ số của mỗi kho đích gồm số đơn tạo, tỷ lệ huỷ, GMV, lead time bình quân và tỷ lệ giao đúng hạn, tính trên các đơn đặt từ ngày 1/7/2017 tới hết tháng 8/2017. Mỗi chỉ số dùng đúng tập đơn trong định nghĩa ở tab Nghiệp vụ.
  \item Truy vấn dưới đây có ba lỗi chồng nhau. Sửa cả ba; giữ nguyên các cột \code{dest_hub_id}, \code{so_don}, \code{ty_le_huy_pct} (phần trăm, làm tròn 2 chữ số), \code{gmv} (USD, làm tròn 0 chữ số), \code{lead_time_tb} (giờ, làm tròn 2 chữ số), \code{ty_le_dung_han_pct} (phần trăm, làm tròn 2 chữ số) và thứ tự sắp.
\end{yeucau}
\begin{erd}
  \erdbang{orders}{0,0}{order_id:PK, origin_hub_id:FK, dest_hub_id:FK, dest_point_id:FK}
  \erdbang{sorting_hubs}{1,0}{hub_id:PK}
  \erdbang{delivery_points}{0,1}{point_id:PK}
  \erdbang{warehouse_daily_ops}{2,0}{hub_id:FK, date}
  \erdbang{staffing_plan}{2,1}{hub_id:FK, date}
  \erdkhoa{Mọi chỉ số tính từ \code{orders}, gom theo \code{dest_hub_id}. \code{sorting_hubs}, \code{delivery_points}, \code{warehouse_daily_ops}, \code{staffing_plan} không cần.}
\end{erd}
\begin{nghiepvu}{Báo cáo chất lượng giao hàng của từng kho đích cho ban giám đốc trên một bộ số nhất quán.}
  \khainiem{ty-le-huy}
  \khainiem{gmv}
  \khainiem{lead-time-don}
  \khainiem{giao-dung-han-don}
  \luuy{Các chỉ số nằm cạnh nhau trên một bảng nhưng không nhất thiết cùng mẫu số. Đơn huỷ vẫn có giá trị ở các cột thời gian và ở cờ đúng hạn.}
\end{nghiepvu}
\goiy{Hai lỗi nằm ở tập đơn của hai chỉ số, lỗi còn lại nằm ở khoảng ngày. Lọc đơn huỷ ngay ở \code{WHERE} thì tỷ lệ huỷ còn đúng không?}
\goiy{\code{AVG(CASE WHEN is_cancelled = 0 THEN lead_time_total_h END)}, không thêm \code{ELSE 0}.}
\ghichuthietke{Khó (sửa v4, viết lại từ câu Dễ \code{scm-a2-ty-le-lo-truot-100}, câu cũ trùng ô với \code{scm-j4}). Đáp án theo từ điển khái niệm của bộ: \code{giao-dung-han-don} lấy mẫu số là đơn không huỷ; \code{lead-time-don} tính tới lúc giao thành công nên đơn huỷ không thuộc tập đơn; \code{ty-le-huy} tính trên mọi đơn tạo. Đáp án khớp sách dự án (\code{main.tex} bản 07/11, thư mục \code{SupplyChain/bai-giang/Final/capnhat_0711}): Mục III.3.11, đoạn ``Đơn huỷ'' ghi ``Nếu quên lọc cột này, lead time và tỷ lệ đúng hạn vẫn ra số trông hợp lý nhưng đã lẫn gần một phần mười số đơn không được giao''; bảng tóm tắt Mục III.3.12 ghi ``Bỏ đơn huỷ khi tính lead time và tỷ lệ đúng hạn''; Bảng thuật ngữ: tỷ lệ huỷ là số đơn khách huỷ chia tổng đơn tạo, GMV không tính đơn huỷ. Dashboard B (Mục VI.5.1) giữ thẻ lead time, thẻ on-time trên mọi đơn chỉ để khớp bản gốc, và sách nêu đó là ba thẻ cạnh nhau đọc trên hai tập đơn khác nhau, sửa bằng cách lọc \code{is_cancelled} cho hai thẻ đó. Ghi chú v4 dẫn phần dashboard B như cách tính của sách là dẫn sai; đã đối chiếu sách theo yêu cầu giảng viên (11/11). Starter (1) lấy lead time và (2) tỷ lệ đúng hạn bình quân trên mọi đơn kể cả đơn huỷ; (3) thiếu cận trên, lỗi này chỉ lộ ở hidden-1 (đơn \code{ORD2_01033328} ngày 2017-09-01). Lọc đơn huỷ ở \code{WHERE} làm tỷ lệ huỷ về 0, nên phải dùng hàm gộp có điều kiện. 5 trục: 1 bảng dùng (chưa đạt); 2 tầng (chưa đạt); mỗi chỉ số một tập đơn theo định nghĩa, học viên tự đối chiếu (Khó); 3 bẫy, 1 chỉ lộ ở hidden-1 (Khó); sai định nghĩa chỉ số (Khó) → Khó 3/5, theo bài Khó của Module 4 trong lộ trình (sửa tỷ lệ trên một bảng). ERD 5 bảng chỉ khoá, 4 bảng nhiễu. Bài sai: lọc đơn huỷ ở \code{WHERE}; \code{ELSE 0} trong \code{AVG}; chỉ sửa lead time; ba bài rớt ở public. Bài sai sửa đủ hai chỉ số nhưng thiếu cận trên rớt ở hidden-1.}
\end{cauhoi}

% ================================================================ A3 · write · Dễ
\begin{cauhoi}[ma=scm-a3-voucher-theo-loai-kho, dulieu=scm, domain=marketing,
  hinhthuc=write, cap=easy, dang=C2, kynang={GROUP BY, SUM, JOIN}, bay={mau-so-sai, quen-can-tren},
  tacgia={Agent gói 2}, trangthai=review, bokiemtra={public, hidden-1}, thutu=co]
\tieude{Khách tự áp mã ở kho chính hay kho vệ tinh nhiều hơn?}
\boicanh[Chị Ngọc · Trưởng nhóm marketing]{Chị muốn biết khách ở vùng kho chính hay vùng kho vệ tinh chủ động dùng voucher nhiều hơn. Em tính giúp chị cho tháng 7/2017, chỉ với nhóm khách tự áp mã.}
\begin{yeucau}
  \item Xét các đơn có \code{order_date} trong tháng 7/2017 thuộc nhóm thử nghiệm \code{checkout_group = 'manual'}.
  \item Tính cho kho chính và kho vệ tinh (cột \code{role} của kho đích) số đơn, số đơn có voucher và tỷ lệ đơn dùng voucher, theo tab Nghiệp vụ.
  \item Trả về \code{role}, \code{so_don}, \code{so_don_voucher}, \code{ty_le_voucher_pct} (phần trăm, làm tròn 2 chữ số).
  \item Sắp \code{role} tăng dần.
\end{yeucau}
\begin{erd}
  \erdbang{orders}{0,0}{order_id:PK, order_date*, dest_hub_id:FK, voucher_applied*, checkout_group*}
  \erdbang{sorting_hubs}{1,0}{hub_id:PK, role*, city}
  \erdnoi[kieu=n-1, tren={\code{orders.dest_hub_id = sorting_hubs.hub_id}}]{orders.dest_hub_id}{sorting_hubs.hub_id}
\end{erd}
\begin{nghiepvu}{Theo dõi mức khách chủ động dùng voucher ở vùng kho chính và vùng kho vệ tinh.}
  \khainiem{voucher}
  \luuy{Mẫu số là mọi đơn của nhóm, kể cả đơn không dùng voucher. Kho đích là kho nhận đơn để giao, khác kho xuất.}
\end{nghiepvu}
\goiy{Ở nhóm hệ thống tự động áp mã, đơn nào cũng có voucher. Gộp nhóm đó vào thì tỷ lệ đo điều gì?}
\goiy{\code{voucher_applied} là 0 hoặc 1, nên \code{SUM(voucher_applied)} là số đơn có voucher.}
\ghichuthietke{Dễ: 2 bảng có điều kiện nối. Sửa v3 (10/11): đơn nhóm \code{'auto_apply'} luôn có \code{voucher_applied} = 1 (sách dự án, phần kiểm tra trước thử nghiệm A/B: 100\% so với 4,15\%), nên chỉ số đổi sang tỷ lệ trong nhóm \code{'manual'} (khoảng 3,9\% ở cả hai loại kho). Bài sai: không lọc nhóm (tỷ lệ quanh 50\%, đo tỷ trọng nhóm A/B); nối theo kho xuất; thiếu cận trên lẫn đơn tháng 8; cả ba rớt ở public.}
\end{cauhoi}

% ================================================================ A4 · write · Dễ
\begin{cauhoi}[ma=scm-a4-cua-hang-vang-khach-2016, dulieu=scm, domain=customer,
  hinhthuc=write, cap=easy, dang=C3, kynang={AVG, GROUP BY, HAVING}, bay={where-thay-having},
  tacgia={Agent gói 2}, trangthai=review, bokiemtra={public, hidden-1}, thutu=co]
\tieude{Cửa hàng nào vắng khách năm 2016?}
\boicanh[Anh Khoa · Quản lý vùng]{Anh đang rà các cửa hàng ít khách để tính chuyện đổi giờ mở cửa. Em tìm giúp anh những cửa hàng năm 2016 có bình quân chưa tới 1.000 lượt giao dịch mỗi ngày.}
\begin{yeucau}
  \item Xét các ngày trong năm 2016.
  \item Tính lượt giao dịch bình quân mỗi ngày của từng cửa hàng, trên các ngày có số liệu giao dịch; chỉ giữ cửa hàng có mức bình quân dưới 1.000.
  \item Trả về \code{store_nbr}, \code{city}, \code{gd_tb_ngay} (làm tròn 0 chữ số).
  \item Sắp \code{gd_tb_ngay} tăng dần; bằng nhau thì \code{store_nbr} nhỏ hơn đứng trước.
\end{yeucau}
\begin{erd}
  \erdbang{transactions}{0,0}{date*, store_nbr:FK, transactions*}
  \erdbang{stores}{1,0}{store_nbr:PK, city*, type, cluster}
  \erdnoi[kieu=n-1, tren={\code{transactions.store_nbr = stores.store_nbr}}]{transactions.store_nbr}{stores.store_nbr}
\end{erd}
\begin{nghiepvu}{Chọn cửa hàng cần đổi giờ mở cửa hoặc chạy chương trình tăng lượt giao dịch.}
  \khainiem{luot-giao-dich}
  \luuy{Ngưỡng 1.000 áp cho mức bình quân của cả năm, không áp cho từng ngày.}
\end{nghiepvu}
\goiy{Điều kiện trên giá trị bình quân chỉ kiểm được sau khi gom nhóm.}
\goiy{Lọc nhóm bằng \code{HAVING AVG(t.transactions) < 1000}, không lọc từng ngày ở \code{WHERE}.}
\ghichuthietke{Dễ: 2 bảng có điều kiện nối. Bài sai lọc từng ngày bằng \code{WHERE transactions < 1000} rồi mới lấy bình quân, rớt ở public.}
\end{cauhoi}

% ================================================================ A5 · debug · Dễ
\begin{cauhoi}[ma=scm-a5-so-ma-ban-cua-hang-44, dulieu=scm, domain=sales,
  hinhthuc=debug, cap=easy, dang=A3, kynang={COUNT DISTINCT, GROUP BY}, bay={dem-dong-thay-vi-ma},
  tacgia={Agent gói 2}, trangthai=review, bokiemtra={public, hidden-1}, thutu=co]
\tieude{Số mã hàng có bán ở cửa hàng 44 vượt cả danh mục}
\boicanh[Chị Thảo · Quản lý cửa hàng]{Báo cáo ghi ngành GROCERY I của cửa hàng mình bán gần ba mươi nghìn mã trong tháng 7, trong khi cả chuỗi chỉ có khoảng bốn nghìn mã hàng. Em kiểm tra lại cách đếm nhé.}
\begin{yeucau}
  \item Truy vấn dưới đây cần đếm số mã hàng (SKU) khác nhau có dòng bán ở cửa hàng 44 trong tháng 7/2017, theo ngành hàng.
  \item Sửa lại cho đúng; giữ nguyên cột \code{family}, \code{so_ma} và thứ tự sắp.
\end{yeucau}
\begin{erd}
  \erdbang{sales}{0,0}{date*, store_nbr*, item_nbr:FK*, unit_sales}
  \erdbang{items}{1,0}{item_nbr:PK, family*}
  \erdnoi[kieu=n-1, tren={một mã, nhiều ngày bán}, canhbao]{sales.item_nbr}{items.item_nbr}
\end{erd}
\begin{nghiepvu}{Đo số mã hàng thực sự có bán để quyết định mã nào giữ trên kệ.}
  \khainiem{luong-ban-ngay}
  \luuy{Một mã hàng bán nhiều ngày thì có nhiều dòng trong \code{sales}.}
\end{nghiepvu}
\goiy{Mỗi dòng của \code{sales} là một mã trong một ngày. Cần đếm gì để mỗi mã chỉ tính một lần?}
\ghichuthietke{Dễ: đường nối tô cảnh báo chỉ ra quan hệ một mã nhiều ngày. Starter \code{COUNT(*)} đếm dòng. Bài sai đếm số ngày có bán (\code{COUNT(DISTINCT `date`)}), rớt ở public.}
\end{cauhoi}

% ================================================================ A6 · debug · Cơ bản
\begin{cauhoi}[ma=scm-a6-tong-giao-dich-thang-7-thap, dulieu=scm, domain=customer,
  hinhthuc=debug, cap=basic, dang=A3, kynang={SUM, GROUP BY}, bay={sai-ham-gop},
  tacgia={Agent gói 2}, trangthai=review, bokiemtra={public, hidden-1}, thutu=co]
\tieude{Tổng lượt giao dịch tháng 7 của cửa hàng thấp bất thường}
\boicanh[Bạn Vy · Thực tập sinh]{Em lập bảng 10 cửa hàng có tổng lượt giao dịch tháng 7/2017 cao nhất, mà cửa hàng nào cũng chỉ vài nghìn lượt cả tháng. Anh chị xem giúp em sai ở đâu.}
\begin{yeucau}
  \item Truy vấn dưới đây cần trả về 10 cửa hàng có tổng lượt giao dịch cao nhất từ ngày 1/7/2017 tới hết ngày 31/7/2017, nhưng số ra quá thấp.
  \item Sửa lại; giữ nguyên cột \code{store_nbr}, \code{tong_gd}, thứ tự sắp và số dòng.
\end{yeucau}
\begin{erd}
  \erdbang{transactions}{0,0}{date*=2017-07-15, store_nbr*=44, transactions*=5051}
\end{erd}
\begin{nghiepvu}{Chọn cửa hàng đông khách để ưu tiên mở thêm quầy thanh toán.}
  \khainiem{luot-giao-dich}
  \luuy{Mỗi dòng là lượt giao dịch của một cửa hàng trong một ngày. Tổng của tháng cộng lượt giao dịch của mọi ngày trong tháng.}
\end{nghiepvu}
\goiy{So giá trị mẫu của một ngày trong ERD với số ra ở cột \code{tong_gd}. Hàm gộp nào cộng các ngày lại?}
\ghichuthietke{Cơ bản (sửa v4, viết lại từ câu Dễ \code{scm-a6-dung-han-theo-mien-thap}, câu cũ trùng ô với \code{scm-j9}). Lỗi một chỗ dễ thấy: starter lấy \code{AVG} thay \code{SUM}, nên \code{tong_gd} chỉ bằng lượt của một ngày bình quân; giá trị mẫu trong ERD là manh mối. 5 trục đều ở mức Cơ bản: 1 bảng; 1 tầng; đề ghi rõ cột, khoảng ngày, số dòng; 0 bẫy; không công thức. Bài sai \code{COUNT(transactions)} đếm số ngày có số liệu, rớt ở public. Hidden-1 bỏ một phần mười số dòng giao dịch nên kết quả khác public.}
\end{cauhoi}

% ================================================================ A7 · write · Trung bình
\begin{cauhoi}[ma=scm-a7-don-mien-theo-quy-2016, dulieu=scm, domain=customer,
  hinhthuc=write, cap=medium, dang=A4, kynang={SUM(CASE WHEN), QUARTER, GROUP BY}, bay={count-else-0, tinh-ca-don-huy},
  tacgia={Agent gói 2}, trangthai=review, bokiemtra={public, hidden-1}, thutu=co]
\tieude{Mỗi miền có bao nhiêu đơn trong từng quý 2016?}
\boicanh[Anh Khoa · Quản lý vùng]{Anh cần một bảng gọn để so mùa vụ đơn online của ba miền, mỗi miền một dòng, mỗi quý một cột. Em làm giúp anh cho năm 2016.}
\begin{yeucau}
  \item Xét các đơn không huỷ, đặt trong năm 2016.
  \item Đếm số đơn của từng miền trong mỗi quý.
  \item Trả về \code{region}, \code{quy_1}, \code{quy_2}, \code{quy_3}, \code{quy_4}.
  \item Sắp \code{region} tăng dần.
\end{yeucau}
\begin{erd}
  \erdbang{orders}{0,0}{order_id:PK, order_date*, region*, dest_hub_id:FK, dest_point_id:FK, is_cancelled*}
  \erdbang{sorting_hubs}{1,0}{hub_id:PK, city, role}
  \erdbang{delivery_points}{1,1}{point_id:PK, city}
  \erdnoi[kieu=n-1]{orders.dest_hub_id}{sorting_hubs.hub_id}
  \erdnoi[kieu=n-1]{orders.dest_point_id}{delivery_points.point_id}
\end{erd}
\begin{nghiepvu}{Lập kế hoạch nhân sự giao hàng theo mùa vụ của từng miền.}
  \khainiem{ty-le-huy}
  \luuy{Đơn huỷ không phát sinh việc giao, nên không tính vào số đơn cần phục vụ.}
\end{nghiepvu}
\goiy{Mỗi cột quý là một phép đếm có điều kiện trên cùng một nhóm miền.}
\goiy{\code{SUM(CASE WHEN QUARTER(order_date) = 1 THEN 1 ELSE 0 END)}; lưu ý \code{COUNT} đếm cả số 0.}
\ghichuthietke{Trung bình: 3 bảng chỉ ký hiệu quan hệ, không cần nối; bảng chéo bằng đếm có điều kiện. Bài sai \code{COUNT(CASE … ELSE 0 END)} ra số cả năm ở mọi cột; bài sai tính cả đơn huỷ; cả hai rớt ở public.}
\end{cauhoi}

% ================================================================ A8 · write · Trung bình
\begin{cauhoi}[ma=scm-a8-bien-ln-gop-nganh-hang, dulieu=scm, domain=finance,
  hinhthuc=write, cap=medium, dang=A3, kynang={SUM, GROUP BY, HAVING, JOIN}, bay={trung-binh-ty-le},
  tacgia={Agent gói 2}, trangthai=review, bokiemtra={public, hidden-1}, thutu=co]
\tieude{Ngành hàng lớn nào có biên lợi nhuận gộp thấp?}
\boicanh[Chị Lan · Giám đốc tài chính]{Chị chuẩn bị đàm phán lại giá với nhà cung cấp của các ngành hàng lớn. Em xếp giúp chị các ngành hàng theo biên lợi nhuận gộp tháng 7/2017.}
\begin{yeucau}
  \item Xét kỳ 2017-07.
  \item Tính doanh thu và biên lợi nhuận gộp của từng ngành hàng theo tab Nghiệp vụ; chỉ giữ ngành có doanh thu từ 1.000.000 trở lên.
  \item Trả về \code{family}, \code{doanh_thu} (làm tròn 0 chữ số), \code{bien_ln_gop_pct} (phần trăm, làm tròn 2 chữ số).
  \item Sắp \code{bien_ln_gop_pct} tăng dần; bằng nhau thì \code{family} theo thứ tự chữ cái.
\end{yeucau}
\begin{erd}
  \erdbang{inventory_snapshots}{0,0}{store_nbr, item_nbr:FK, ym*, sold*, inventory_value}
  \erdbang{item_costs}{1,0}{item_nbr:PK, family*, unit_price*, unit_cost*}
  \erdbang{items}{2,0}{item_nbr:PK, family, perishable}
  \erdnoi[kieu=n-1]{inventory_snapshots.item_nbr}{item_costs.item_nbr}
  \erdnoi[kieu=1-1]{item_costs.item_nbr}{items.item_nbr}
\end{erd}
\begin{nghiepvu}{Chọn ngành hàng cần đàm phán lại giá mua hoặc điều chỉnh giá bán.}
  \khainiem{bien-ln-gop}
  \khainiem{luong-ban-va-nhu-cau}
  \luuy{Doanh thu của ngành là tổng doanh thu các mã hàng thuộc ngành trong kỳ.}
\end{nghiepvu}
\goiy{Biên của ngành bằng tổng lợi nhuận gộp chia tổng doanh thu. Trung bình biên của từng dòng cho mỗi dòng trọng số như nhau.}
\goiy{Ngưỡng doanh thu áp cho tổng của ngành, nên lọc sau khi gom nhóm.}
\ghichuthietke{Trung bình: 3 bảng chỉ ký hiệu quan hệ; \code{items} không cần vì \code{item_costs} đã có \code{family}. Bài sai lấy trung bình biên từng dòng; bài sai lọc doanh thu từng dòng bằng WHERE (ra 0 dòng); cả hai rớt ở public.}
\end{cauhoi}

% ================================================================ A9 · write · Khó
\begin{cauhoi}[ma=scm-a9-pheu-don-mua-theo-dc-2016, dulieu=scm, domain=operations,
  hinhthuc=write, cap=hard, dang=D1, kynang={COUNT(CASE WHEN), JOIN, GROUP BY, CTE, LEAST}, bay={buoc-khong-long-nhau, sai-mau-so-buoc, count-dem-ca-false, can-ngay-hen},
  tacgia={Agent gói 2}, trangthai=review, bokiemtra={public, hidden-1}, thutu=co]
\tieude{Đơn mua năm 2016 của mỗi DC hụt nhiều nhất ở bước nào?}
\boicanh[Chị Hằng · Trưởng phòng mua hàng]{Chị chuẩn bị làm việc lại với nhà cung cấp của từng trung tâm phân phối (DC). Em xem giúp chị năm 2016 đơn mua của mỗi DC hụt nhiều nhất ở khâu giao đúng hạn, giao đủ hàng hay kiểm định.}
\begin{yeucau}
  \item Xét các đơn mua đặt trong năm 2016, xếp vào DC mà nhà cung cấp giao hàng về.
  \item Dựng phễu đơn mua theo tab Nghiệp vụ và tính tỷ lệ qua từng bước cho mỗi DC; bước hụt nhiều nhất là bước có tỷ lệ qua thấp nhất, bằng nhau thì lấy bước đứng trước.
  \item Trả về \code{warehouse_id}, \code{so_don_dat}, \code{ty_le_dung_han_pct}, \code{ty_le_du_hang_pct}, \code{ty_le_dat_pct} (ba cột tỷ lệ là phần trăm, làm tròn 2 chữ số), \code{buoc_hut_nhat} (một trong ba giá trị \code{'dung_han'}, \code{'du_hang'}, \code{'dat_kiem_dinh'}).
  \item Sắp \code{warehouse_id} tăng dần.
\end{yeucau}
\begin{erd}
  \erdbang{purchase_orders}{0,0}{po_id:PK, supplier_id:FK, item_nbr:FK}
  \erdbang{suppliers}{1,0}{supplier_id:PK, warehouse_id:FK}
  \erdbang{warehouses}{2,0}{warehouse_id:PK}
  \erdbang{items}{0,1}{item_nbr:PK}
  \erdbang{shipments}{2,1}{shipment_id:PK, warehouse_id:FK, store_nbr:FK}
  \erdkhoa{\code{purchase_orders.supplier_id} nối \code{suppliers}; DC của đơn lấy ở \code{suppliers.warehouse_id}. \code{warehouses}, \code{items}, \code{shipments} không cần.}
\end{erd}
\begin{nghiepvu}{Chọn DC và khâu cần làm việc với nhà cung cấp trước: giữ đúng ngày hẹn, giao đủ hàng hay siết chất lượng.}
  \quytrinh{Vòng đời đơn mua}{đặt hàng, hàng về, kiểm định}{chưa về}
  \khainiem{pheu-don-mua}
  \khainiem{otd-ncc}
  \khainiem{lo-dat-lo-truot}
  \luuy{Một đơn đạt kiểm định nhưng về trễ hạn vẫn dừng ở bước giao đúng hạn, không được tính ở các bước sau.}
\end{nghiepvu}
\goiy{Mỗi bước là điều kiện của bước trước cộng thêm một điều kiện mới. Tỷ lệ của một bước chia cho số đơn của bước liền trước, không chia cho số đơn đặt.}
\goiy{Số đơn qua bước cuối là \code{COUNT(CASE WHEN po.received_date <= po.promised_date AND po.qty_received = po.qty_ordered AND po.inspection_result = 'Pass' THEN 1 END)}.}
\ghichuthietke{Khó (sửa v4, viết lại từ câu Trung bình \code{scm-a9-ncc-hang-loi-vuot-nguong}; bẫy trung bình tỷ lệ của câu cũ vẫn có ở a8). Dạng D1 (phễu), bộ chưa có câu nào. Đơn đặt năm 2016 đều đã về nên bước đặt hàng không lẫn đơn chưa tới hạn. 5 trục: 2 bảng dùng (chưa đạt); 4 tầng, cờ từng bước, đếm có điều kiện theo DC, tỷ lệ trên bước trước, chọn bước thấp nhất (Khó); đề nói bằng ngôn ngữ kinh doanh, học viên tự dựng điều kiện lồng nhau (Khó); 4 bẫy ở public, chưa có bẫy chỉ lộ ở hidden-1 (hidden-1 bỏ một phần mười đơn mua nên kết quả khác public); chỉ số nhiều điều kiện (Khó) → Khó 3/5. ERD 5 bảng chỉ khoá, 3 bảng nhiễu. Bài sai: chia mọi tỷ lệ cho số đơn đặt; bước sau không giữ điều kiện của bước trước; \code{COUNT(điều kiện)} đếm cả 0; về đúng hạn dùng \code{<} thay \code{<=}; cả bốn rớt ở public.}
\end{cauhoi}


% ################################################################ CASE STUDY

% ================================================================ C3 · case · Cơ bản
\begin{cauhoi}[ma=scm-c3-kho-giao-tre, dulieu=scm, domain=operations,
  hinhthuc=case, cap=basic, diem=10, tacgia={Agent gói 2}, trangthai=review, bokiemtra={public, hidden-1}, thutu=co]
\tieude{Kho đích nào cần ưu tiên khắc phục giao trễ?}
\boicanh[Anh Phúc · Trưởng kho Hà Nội]{Ban giám đốc hỏi quý 1/2016 kho nào giao trễ đáng lo nhất để ưu tiên hỗ trợ. Em xem số liệu rồi cho anh kết luận.}
\begin{erd}
  \erdbang{orders}{0,0}{order_id:PK, order_date*=2016-02-20, dest_hub_id*='HUB_HUE', on_time*=0, is_cancelled*=0}
\end{erd}
\begin{nghiepvu}{Chọn kho cần ưu tiên hỗ trợ nhân sự và quy trình giao hàng.}
  \khainiem{giao-dung-han-don}
  \luuy{Kho đích là kho nhận đơn để giao tới khách; mỗi đơn chỉ có một kho đích.}
\end{nghiepvu}
\begin{buoc}[ma=b1, diem=5, hinhthuc=write, bang={orders}]{Tỷ lệ trễ theo kho đích}
  \begin{yeucau}
    \item Xét đơn không huỷ, đặt từ ngày 1/1/2016 tới hết ngày 31/3/2016. Đơn trễ là đơn có \code{on_time} bằng 0.
    \item Trả về \code{dest_hub_id}, \code{so_don}, \code{so_don_tre}, \code{ty_le_tre_pct} (phần trăm, làm tròn 2 chữ số).
    \item Sắp \code{ty_le_tre_pct} giảm dần; bằng nhau thì \code{dest_hub_id} theo thứ tự chữ cái.
  \end{yeucau}
  \goiy{Đếm đơn trễ bằng \code{COUNT(CASE WHEN on_time = 0 THEN 1 END)}.}
\end{buoc}
\begin{buoc}[ma=b2, diem=5, hinhthuc=ketluan, bang={orders}]{Kết luận}
  \begin{luachon}
    \phuongan{\code{DC_HANOI} cần được ưu tiên hỗ trợ trước, vì số đơn trễ của kho này nhiều hơn hẳn năm kho còn lại.}
      {\code{DC_HANOI} nhận hơn bảy phần mười số đơn, nên số đơn trễ lớn dù tỷ lệ trễ chỉ đứng thứ ba.}
    \phuongan{\code{HUB_PHANTHIET} có tỷ lệ trễ cao nhất trong sáu kho, nên có lẽ cần ưu tiên hỗ trợ kho này trước.}
      {\code{HUB_PHANTHIET} chỉ có khoảng 30 đơn trễ; bớt bốn đơn trễ là kho này đã xuống dưới kho thứ hai.}
    \phuongan*{Ba kho đứng đầu có tỷ lệ trễ sát nhau, nên số liệu quý 1 chưa đủ để chọn riêng một kho ưu tiên.}
      {Đúng: ba kho đầu chênh nhau chưa tới 0,6 điểm phần trăm; tỷ lệ cao nhất của \code{HUB_PHANTHIET} chỉ dựa trên chưa tới 750 đơn.}
  \end{luachon}
\end{buoc}
\ghichuthietke{Case Cơ bản: GROUP BY một cột trên một bảng, theo case mẫu c1 (QUYET-DINH-CHUNG 1.2); 1 bảng có giá trị mẫu, 1 bước SQL, 1 bước kết luận. Phương án sai: lấy số đơn trễ tuyệt đối thay cho tỷ lệ; coi tỷ lệ cao nhất trên ít đơn là khác biệt thật. Sửa V4 (11/11): thu kỳ từ cả năm 2016 về quý 1/2016 để truy vấn chạy nhanh (gom cả năm theo \code{dest_hub_id} làm MySQL quét theo chỉ mục kho, 8--22 giây; quý 1 dùng chỉ mục ngày, khoảng 1,3 giây); bỏ phương án thứ tư cho đủ ba phương án; thêm từ rào vào một phương án sai để phương án đúng không lộ. Số liệu quý 1: \code{HUB_PHANTHIET} 4,52\% (33/730 đơn), \code{HUB_THAIBINH} 3,98\%, \code{DC_HANOI} 3,94\% với 3.016 đơn trễ; hidden-1: 4,41\% (29/658), 3,94\%, 3,93\%. Kết luận đúng ở cả hai biến thể. Bài sai bước 1 sắp theo số đơn trễ, rớt ở public.}
\end{cauhoi}

% ================================================================ C4 · case · Dễ
\begin{cauhoi}[ma=scm-c4-khuyen-mai-cua-hang-44, dulieu=scm, domain=marketing,
  hinhthuc=case, cap=easy, diem=20, tacgia={Agent gói 2}, trangthai=review, bokiemtra={public, hidden-1}, thutu=co]
\tieude{Hàng khuyến mãi ở cửa hàng 44 bán cao hơn bao nhiêu?}
\boicanh[Chị Ngọc · Trưởng nhóm marketing]{Chị đang chia ngân sách khuyến mãi quý sau cho cửa hàng 44. Em xem số liệu tháng 7/2017 rồi cho chị biết nên ưu tiên khuyến mãi cho ngành hàng nào.}
\begin{erd}
  \erdbang{sales}{0,0}{date*, store_nbr*, item_nbr:FK, unit_sales*, onpromotion*}
  \erdbang{items}{1,0}{item_nbr:PK, family*, class, perishable}
  \erdnoi[kieu=n-1, tren={\code{sales.item_nbr = items.item_nbr}}]{sales.item_nbr}{items.item_nbr}
\end{erd}
\begin{nghiepvu}{Chia ngân sách khuyến mãi giữa các ngành hàng của một cửa hàng.}
  \khainiem{khuyen-mai}
  \khainiem{luong-ban-ngay}
  \luuy{Lượng bán bình quân mỗi dòng là tổng lượng bán chia số dòng bán; mỗi dòng là một mã hàng trong một ngày.}
\end{nghiepvu}
\begin{buoc}[ma=b1, diem=8, hinhthuc=write, bang={sales}]{Lượng bán bình quân mỗi dòng khi có và không có khuyến mãi}
  \begin{yeucau}
    \item Xét các dòng bán của cửa hàng 44 trong tháng 7/2017.
    \item Trả về \code{onpromotion}, \code{so_dong}, \code{luong_ban_tb} (lượng bán bình quân mỗi dòng, làm tròn 2 chữ số); sắp \code{onpromotion} tăng dần.
  \end{yeucau}
  \goiy{Lọc ngày có cả cận dưới và cận trên, rồi gom theo \code{onpromotion}.}
\end{buoc}
\begin{buoc}[ma=b2, diem=6, hinhthuc=write, bang={sales, items}]{Mức tăng theo ngành hàng}
  \begin{yeucau}
    \item Cùng phạm vi bước 1, chỉ xét sáu ngành \code{'BEVERAGES'}, \code{'CLEANING'}, \code{'DAIRY'}, \code{'DELI'}, \code{'GROCERY I'}, \code{'PRODUCE'}.
    \item Trả về \code{family}, \code{tb_co_km}, \code{tb_khong_km} (lượng bán bình quân mỗi dòng khi có và khi không khuyến mãi, làm tròn 2 chữ số); sắp \code{family} theo thứ tự chữ cái.
  \end{yeucau}
  \goiy{\code{AVG(CASE WHEN onpromotion = 1 THEN unit_sales END)}; không thêm \code{ELSE 0}.}
\end{buoc}
\begin{buoc}[ma=b3, diem=6, hinhthuc=ketluan, bang={sales, items}]{Kết luận}
  \begin{luachon}
    \phuongan{Khuyến mãi làm lượng bán mỗi dòng tăng khoảng gấp đôi ở mọi ngành hàng, nên chia đều ngân sách cho các ngành.}
      {Mức gấp đôi là số chung của bước 1; bước 2 cho thấy mức tăng rất khác nhau giữa các ngành.}
    \phuongan*{Chênh lệch khi có khuyến mãi khác nhau nhiều giữa các ngành, lớn ở BEVERAGES, CLEANING và gần như không có ở DAIRY, nên thử ưu tiên hai ngành đầu.}
      {Đúng: bước 2 cho thấy chênh lệch lớn ở BEVERAGES, CLEANING và rất nhỏ ở DAIRY. Mã được chọn khuyến mãi có thể vốn bán chạy.}
    \phuongan{Khuyến mãi ít tác dụng vì số dòng có khuyến mãi chỉ chiếm phần nhỏ trong tổng số dòng bán của cửa hàng.}
      {Tỷ trọng dòng có khuyến mãi cho biết khuyến mãi phủ rộng tới đâu, không cho biết mức tăng lượng bán.}
  \end{luachon}
\end{buoc}
\ghichuthietke{Case Dễ: ERD cố định 2 bảng có điều kiện nối; bước 1 chỉ dùng \code{sales}, bước 2 tô sáng thêm \code{items}. Bài sai bước 1 thiếu cận trên rớt ở hidden-1 (dòng cửa hàng 44 ngày 2017-08-01); bài sai bước 2 \code{ELSE 0} trong AVG rớt ở public.}
\end{cauhoi}

\end{document}
